%%%%%%%%%%%%%%%%%%%%%%%%%%%%%%%%%%%%%%%%%%%%%%%%%%%%%%%%%%%%%%%
%
% Welcome to Overleaf --- just edit your LaTeX on the left,
% and we'll compile it for you on the right. If you open the
% 'Share' menu, you can invite other users to edit at the same
% time. See www.overleaf.com/learn for more info. Enjoy!
%
%%%%%%%%%%%%%%%%%%%%%%%%%%%%%%%%%%%%%%%%%%%%%%%%%%%%%%%%%%%%%%%
\documentclass[17pt, a0paper, landscape,  margin=0mm, innermargin=10mm, blockverticalspace=3mm, colspace=3mm, subcolspace=3mm]{tikzposter}
\usepackage{graphicx}
\usepackage{authblk}
\usepackage{color}
% \usepackage[all]{xy}

\usetheme{board}

%% Elena's favorite green (thanks, Fernando!)
\definecolor{ForestGreen}{RGB}{34,139,34}
\definecolor{BlueViolet}{RGB}{138,43,226}
\definecolor{Coquelicot}{RGB}{255, 56, 0}
\definecolor{Teal}{RGB}{2,132,130}
%Uncomment this if you want to show work-in-progress comments
\newcommand{\comment}[1]{{\bf \tt  {#1}}}
% Uncomment this if you don't want to show comments
%\newcommand{\comment}[1]{}
% for color names see https://www.overleaf.com/learn/latex/Using_colors_in_LaTeX
\newcommand{\emcomment}[1]{\textcolor{ForestGreen}{\comment{Elena: {#1}}}}

\newcommand{\jscomment}[1]{\textcolor{magenta}{\comment{James: {#1}}}}


\title{Developing a framework for text-based game development for beginners in Clojure}
\author[1]{James Swearingen}
\author[1]{Elena Machkasova (adviser)}
\date{\today}
\affil[1]{Division of Science and Mathematics, University of Minnesota, Morris}
\usetheme{Envelope}
\usecolorstyle{Australia}

\begin{document}

\makeatletter
\def\maketitle{\AB@maketitle}
\makeatother

\maketitle

\block{Abstract}{
Clojure is a relatively new functional Lisp-like programming language that is used in industry and is supported by a community of developers. It is also a promising language for a beginner classroom due to its simple syntax, focus on small functions and on recursive algorithms, and data that’s immutable by default. However, there are gaps that make its adoption challenging. Debugging in Clojure can be unintuitive, because Clojure errors are exceptions in the underlying Java language. They frequently contain terminology intended for experienced developers, leaving out important context.

Past work at Morris included the modifications to the error messages to make them understandable by beginners, but in order to perform usability testing with the target audience the error messages need to be incorporated into a simple development environment. This project is a continuation of an ongoing research, focusing on developing such an environment.

% My work to be presented at MICS is focused on developing a framework for students to create simple text-based interactive games. The framework handles modifications to the game state while allowing students to write functions to describe the game state transformations without directly using mutability. Additionally, I am working on a set of simple functions that allow students to test their code without using sophisticated testing engines. This environment is inspired by a similar environment in the Racket programming language successfully used in introductory programming classes, although at this point we are not including graphical features used in Racket.

% My poster will present examples of uses of the environment. It will also discuss implementation challenges and their solutions. 
%\emcomment{Needs to be shortened and changed to match this year's work}
}
\begin{columns}
\column{0.3}
        \block{Overview of Functional Programming and Clojure}{
%\tkcomment{Tristan's section}
    \begin{itemize}
%	\item \emcomment{Includes reasons to teach FP to beginners}
	\item Functional programming has a rich history of being introduced early in programming education [4].
	\item It emphasizes breaking problems into smaller, manageable, easy-to-combine pieces.
%promotes a more declarative programming style. Which is where you tell it what you want to happen rather than giving it a list of steps to follow.
	\item It builds a strong foundation for students in subjects like recursion, higher-order functions, immutable data, and problem decomposition.
    \item Clojure is a functional language
%Lisp-based language 
that encourages these principles and is a great tool for introducing students to functional programming.
    \item It has a simple, minimal syntax with prefix notation in statements and statements are surrounded by parentheses: \texttt{(+ 2 3)} is the way to write $2 + 3$.
    \item It uses an interactive Read-Evaluate-Print-Loop (REPL) environment: the user types some code, the system evaluates it and prints the answer.
    \end{itemize}
}

 \block{Clojure Error Messages}{
%\tkcomment{Tristan's section}
    \begin{itemize}
    \item Error messages are understandable to experienced developers, but have terminology that does not make sense to beginner programmers.
    \item Clojure code is compiled into a more common language, Java. It is a more complex language whose data types and features are not known to beginner programmers. 
   \item By the time Clojure code is executed, it is Java code, so Clojure errors appear as Java errors.
%\emcomment{The target audience doesn't know the difference.}
    \item Error messages include Java error types and a Java-phrased cause, as well as a Clojure location.
    \end{itemize}
}

\block{Our Project, Babel}{
    \begin{itemize}
        \item One focus of our research is development of a tool called Babel to modify Clojure error messages.
        \item Many error messages are directly replaced in the REPL with simplified, beginner-friendly versions.
        \item Errors on function calls provide more detailed information about expected arguments: their types, number, etc.
        \item Babel categorizes error data via dictionary lookup and specifications on Clojure function requirements to achieve this. 
        \item We developed and designed these specifications to capture as much information for the end user as possible, which is absent in default Clojure messages.
    \end{itemize}
}


\block{Examples}{
%\jwcomment{John's section.}
%\emcomment{(even? "two"), (even? 2 3)}
    Correct call: \\
    \texttt{user=>(even? 2)} \\
    \texttt{true} \\
\jscomment{This is probably much too wordy for what you were suggesting, and it definitely doesn't fit on the left column here without also moving something else over or cutting a lot from the other sections.}

    Default Clojure Error Messages in REPL: \\
    \texttt{user=>(even? "two")} \\
    \texttt{Execution error (IllegalArgumentException) at user/eval2044 (REPL:1).}\\
    \texttt{Argument must be an integer: two}
The message gives the error as a Java exception, which isn't helpful to a beginner programmer. Also, by default, strings are printed without quotation marks into the console, making it unclear what the failing argument is.

    \texttt{user=>(even? 2 3)} \\
    \texttt{Execution error (ArityException) at user/eval2046 (REPL:1).} \\ 
     \texttt{Wrong number of args (2) passed to: clojure.core/even?} \\
As in the previous error, the exact call to the failing function is not stated clearly. The only reference we get that \texttt{(even? 2 3)} was the exact call that failed was the \texttt{clojure.core/even?}. This message may be ambiguous to a beginner if they called this function in multiple places.

    Default Babel Error Messages in REPL: \\
    \texttt{user=>(even? "two")}\\
    \texttt{The first argument of (even? "two") was expected to be a number but is a string "two" instead.}\\
    \texttt{In Clojure interactive session on line 1.}
This message replaces the confusing programming jargon with a simple English sentence that explains the error and what it was. Babel performs similar replacements with almost all Java and Clojure exceptions. In addition, the string is given quotation marks, making it clear that it was a string.

    \texttt{user=>(even? 2 3)} \\
    \texttt{Wrong number of arguments in (even? 2 3): the function even? expects one argument but was given two arguments.}\\
    \texttt{In Clojure interactive session on line 1.} \\
The specific function call is stated in the error message, and the statement about "wrong number of args" is instead phrased in clear English.

}

\column{0.4}
 \block{Goals for This Year}{
    \begin{itemize}
        \item Presently, we are working with a tool called Morse to provide an interactive user interface for Babel.
        \item Morse is a data visualization tool designed for professional development and debugging [2].
        \item In the past, we succeeded in integrating Morse with Babel. Now, we are modifying information flow between the tools to provide Morse with more capabilities.
        \item We are aiming to have Morse stylize text based on good design practices for error messages [3].
    \end{itemize}
    }

\column{0.3}
\block{Progress This Year}{
%\jwcomment{John's section.}
    \begin{itemize}
        \item We improved internal tooling to allow us to easily select different types of error messages and test how they are presented.
        \item We implemented a skeleton framework for data flow to the interactive viewer.
        \item We created a system for labeling parts of an error message to be able to use different formats for different parts.
        \item We created a simple viewer to format error messages based on labeled data.
    \end{itemize}
}

\block{Future Work}{
%\jscomment{Jaydon's section.}
    \begin{itemize}
        \item We would like to develop more Morse viewers for other information provided in error messages, such as stack traces and Java errors, to allow users to gather information about different parts of an error through Morse.
        \item We are looking to conduct usability studies on displaying error messages via Morse in the future.
        \item We also want to explore integrating Morse viewers into IDEs such as Visual Studio Code for workflow refinements.
        \item We are considering adding support for defining terms in messages via hover text, as well as providing more access to resources like documentation.
    \end{itemize}
}

\block{Acknowledgments}{
 Thanks to Morris Academic Partnership (MAP) and the Undergraduate Research Opportunity (UROP) for funding this work. \\
Thanks to Joe Lane (Nubank) for guidance on tools for this project. 
}

\block{References}{
    \begin{itemize}
        \item {[1]} \emph{clojure.org}
        \item {[2]} Morse, Nubank \emph{https://github.com/nubank/morse}
	\item {[3]} Tao Dong, Kandarp Khandwala,  \textit{The Impact of ``Cosmetic'' Changes on the Usability of Error Messages}, 2019 CHI Conference on Human Factors in Computing Systems.
	\item {[4]} Matthias Felleisen, Robert Bruce Findler, Matthew Flatt, and Shriram Krishnamurthi \textit{The Structure and Interpretation of the Computer Science Curriculum},  Journal of Functional Programming, 2004.
    \end{itemize}
}

%\block{}{
%\begin{tikzfigure}
%    \includegraphics[width=0.55\linewidth]{figures/UMM logo.png}
%\end{tikzfigure}
%}

\end{columns}


\end{document}
